\documentclass[11pt]{article}

\usepackage{kotex}
\usepackage{geometry}
\usepackage{setspace}
\usepackage{parskip}

\geometry{a4paper,margin=24mm}
\setstretch{1.18}
\setmainhangulfont{AppleMyungjo}
\setsanshangulfont{Apple SD Gothic Neo}

\title{\textbf{프로젝트 기반 자기소개서 초안}\\와이즈와이어즈 지원용}
\author{이준영}
\date{2026년 3월}

\begin{document}

\maketitle

\section*{1. 요구사항부터 시스템 시험까지 닫아 본 검증 경험}
현대모비스 부트캠프 PBL에서 저는 구간 정보, 긴급차량 접근, 객체 위험을 하나의 경고 정책으로 통합하는 프로젝트를 수행했습니다. 이 프로젝트는 차량 전체 기준선 위에서 다수 ECU, CAN/Ethernet 네트워크, V2X·ADAS·CGW 논리를 함께 맞춰야 했기 때문에 기능 구현과 검증 기준을 동시에 정리해야 하는 과제였습니다. 저는 그 과정에서 아키텍처와 통신 구조를 정리하고, 핵심 CAPL 로직과 시험 구조를 함께 맞추는 역할을 맡았습니다.

처음에는 기능 정의와 코드만 맞으면 될 것이라고 생각했습니다. 하지만 실제로는 요구사항, 기능 정의, 네트워크 흐름, 시스템 변수, 코드, 시험 기준이 따로 움직이면 구현은 되어도 검증이 닫히지 않았고, 검증이 되어도 그 근거를 설명하기 어려웠습니다. 그래서 ``무엇을 왜 이렇게 만들었는지''와 ``어떻게 검증할 것인지''를 처음부터 하나의 기준선 위에서 맞추는 일이 중요하다는 점을 체감했습니다.

그 결과 이번 프로젝트에서는 요구사항부터 기능 정의, 통신, 변수, 코드, 단위시험·통합시험·시스템시험까지 이어지는 추적 구조를 정리할 수 있었습니다. 각 단계에서 무엇을 확인해야 하는지 기준을 분리해 두었기 때문에, 기능 로직 검증과 차량 기준선 검증을 구분해서 볼 수 있었습니다. 이 경험을 통해 전장 시스템은 기능 구현만이 아니라, 구현과 검증이 동시에 성립하는 구조로 설계해야 한다는 점을 배웠습니다.

\section*{2. 복잡한 시스템을 설명 가능한 결과로 바꾸는 역량}
이번 프로젝트에서 어려웠던 점은 기능 자체보다, 복잡한 시스템을 팀원과 리뷰어가 같은 기준으로 이해할 수 있게 정리하는 일이었습니다. 구간 경고, V2X, ADAS, 감속 보조, Fail-Safe가 동시에 연결되어 있었기 때문에 작은 표현 차이도 전체 구조를 흔들 수 있었습니다.

저는 이 문제를 코드 수정만으로 풀기보다, 먼저 의미를 정리하는 방식으로 접근했습니다. 어떤 ECU가 어떤 판단을 소유하는지, 어떤 메시지가 어떤 상태로 이어지는지, 시험에서 무엇을 확인해야 하는지를 먼저 정리했고, 그 기준에 맞춰 문서와 구현을 함께 맞췄습니다. 이후에는 UT·IT·ST 자산을 기준으로 시험 실행 결과를 XML 기반으로 정리하고, 엑셀 산출물과 closeout 흐름까지 연결해 반복 작업을 줄이는 방향으로 운영 자동화도 정리했습니다. 이 과정에서 제가 특히 중요하게 본 것은, 실행 결과가 단순 기록으로 끝나는 것이 아니라 ``어떤 요구사항과 기능을 확인한 결과인지''까지 다시 설명될 수 있어야 한다는 점이었습니다.

이 과정을 통해 시스템 프로젝트에서는 문제를 해결하는 것과 함께, 그 결과를 다시 설명 가능하게 정리하는 일이 중요하다는 점을 배웠습니다. 와이즈와이어즈에서도 이 경험을 바탕으로 고객과 개발팀이 함께 이해할 수 있는 검증 결과와 품질 근거를 정리하는 데 기여하고 싶습니다.

\section*{3. 문서와 품질 기준을 실제 구현과 연결한 경험}
프로젝트를 하면서 가장 크게 느낀 점 중 하나는, 문서는 나중에 정리하는 결과물이 아니라 처음부터 구현을 이끄는 기준이어야 한다는 점이었습니다. 실제로 프로젝트 초반에는 기능 설명은 있는데 시험 기준이 약하거나, 코드가 있는데 그 근거가 문서로 남지 않는 순간들이 반복되었습니다. 그럴 때마다 팀은 같은 내용을 다시 논의해야 했고, 구현보다 정리가 더 늦어지는 문제가 생겼습니다.

저는 이런 비효율을 줄이기 위해 문서를 형식적인 산출물이 아니라 ``팀이 같은 말을 쓰게 만드는 기준''으로 다루려고 했습니다. 요구사항은 무엇을 만들지 정의하고, 기능 정의는 어떻게 동작해야 하는지 설명하고, 통신과 변수 문서는 그 동작을 실제 구현과 연결하고, 시험 문서는 그 동작을 무엇으로 확인할지 보여주도록 역할을 나누었습니다. 특히 시험 문서를 작성할 때는 단순히 시나리오를 적는 것이 아니라, 어떤 조건에서 어떤 결과를 기대하고 어떤 근거로 확인할지를 함께 남기려고 했습니다. 덕분에 팀원 간 해석 차이가 줄었고, 구현과 시험을 연결할 때도 훨씬 수월해졌습니다.

이 경험을 통해 품질은 테스트의 양만으로 결정되는 것이 아니라, 처음부터 기준이 흔들리지 않도록 관리하고 결과를 다시 확인할 수 있게 남기는 일과도 연결되어 있다는 점을 배웠습니다. 와이즈와이어즈에서도 문서와 테스트를 따로 보지 않고, 품질을 설명 가능한 형태로 유지하는 데 기여하고 싶습니다.

\section*{4. 팀을 끌고 가는 리더십보다, 기준을 세우는 리더십}
이번 프로젝트에서 저는 팀장 역할을 맡아 아키텍처, 통신 구조, 전체 방향을 조율하는 일을 했습니다. 프로젝트 범위가 넓다 보니 팀원마다 보는 관점이 달랐고, 구현해야 할 기능도 많았습니다. 이럴 때 제 판단만 앞세우기보다, 팀이 공통으로 따를 수 있는 기준을 먼저 정리하려고 했습니다.

대신 제가 하려고 했던 일은 팀이 같은 기준으로 움직이게 만드는 것이었습니다. 어떤 기능을 먼저 봐야 하는지, 어떤 문서가 기준이 되는지, 어떤 시험이 실제로 중요한지를 계속 정리했고, 회의 때도 ``누가 맞는가''보다 ``무엇을 기준으로 판단할 것인가''를 먼저 맞추려고 했습니다. 그 결과 의견 차이가 있더라도 다시 기준선으로 돌아와 논의할 수 있었고, 프로젝트가 커져도 팀이 쉽게 흔들리지 않았습니다.

이번 경험을 통해 리더십은 사람을 앞에서 끌고 가는 역할보다, 팀이 끝까지 같은 방향을 볼 수 있게 만드는 역할에 가깝다는 점을 배웠습니다. 실제 업무에서도 팀이 같은 품질 기준 아래 안정적으로 움직일 수 있도록 기여하는 사람이 되고 싶습니다.

\section*{5. 와이즈와이어즈에서 만들고 싶은 검증의 역할}
제가 와이즈와이어즈에 지원하는 이유는 테스트를 ``오류를 찾는 과정''에만 머무르지 않고, 시스템이 왜 신뢰할 수 있는지를 설명하는 과정이라고 생각하기 때문입니다. 이번 프로젝트에서도 기능이 동작하는 장면 자체보다, 그 기능이 어떤 기준으로 설계되었고 어떤 시험으로 확인되었는지를 설명할 수 있게 되는 과정이 더 중요했습니다.

와이즈와이어즈에서도 저는 단순히 테스트 케이스를 수행하는 역할에 머무르지 않고, 고객과 개발팀이 모두 납득할 수 있는 검증 근거를 만드는 QA가 되고 싶습니다. 특히 모빌리티와 전장 소프트웨어처럼 작은 오류가 큰 문제로 이어질 수 있는 환경에서는 테스트의 양뿐 아니라, 어떤 기준으로 검증했고 어떤 근거로 결과를 판단했는지가 중요하다고 생각합니다.

입사 후에는 현업의 검증 프로세스를 빠르게 익히고, 프로젝트에서 쌓은 전주기 관점과 추적성 중심의 사고방식을 바탕으로, ``무엇을 검증했고 왜 이 결과를 신뢰할 수 있는지''를 분명히 설명할 수 있는 QA 엔지니어로 성장하겠습니다.

\end{document}
